% Generated from locale/tr/content/set-theory/cardinals/classing.tex; do not hand-edit.


\olfileid[tr]{sth}{cardinals}{classing}
\olsection{Sonlu, \usetoken{S}{enumerable}, \usetoken{S}{nonenumerable}}

Kardinal sayılarla tanıştığımıza göre kardinal sayıların farklı türlerini,
özellikle sonlu, !!{enumerable} ve !!{nonenumerable} kardinal sayıları
tartışmak için biraz zaman ayırmak yararlıdır.

İlk iki sonucumuz, sonlu kardinal sayıların tam olarak
\olref[z][infinity-again]{defnomega} içinde \emph{doğal sayılarımız}
olarak tanımladığımız sonlu sıra sayıları olduğunu gerektirir:

\begin{prop}\ollabel{finitecardisoequal}
$n,m\in\omega$ olsun. $n=m$ ancak ve ancak $\cardeq{n}{m}$.
\end{prop}

\begin{proof}
\emph{Soldan sağa} yön açıktır. \emph{Sağdan sola} yön için $n\neq m$
olduğu hâlde $\cardeq{n}{m}$ olduğunu varsayın. Üçlü Karşılaştırma gereği
ya $n\in m$ ya da $m\in n$ olur; genelliği kaybetmeden $n\in m$ varsayın.
O zaman $n\subsetneq m$ ve bir !!a{bijection} $f\colon m\to n$ vardır;
dolayısıyla $m$ Dedekind-sonsuzdur. Bu,
\olref[z][infinity-again]{naturalnumbersarentinfinite} ile çelişir.
\end{proof}

\begin{cor}\ollabel{naturalsarecardinals}
$n\in\omega$ ise $n$ bir kardinal sayıdır.
\end{cor}

\begin{proof}
Doğrudan çıkar.
\end{proof}
\noindent
Bir kardinal sayıyı ``sonlu'' ya da ``sonsuz'' diye betimlemenin çeşitli
makul anlamlarının çakıştığı da çıkar:
\begin{thm}\ollabel{generalinfinitycharacter}
Her $A$ kümesi için aşağıdakiler denktir:
\begin{enumerate}
	\item\ollabel{card:notinomega} $\card{A}\notin\omega$; yani $A$ bir
	doğal sayı değildir;
	\item\ollabel{card:omegaplus} $\omega\leq\card{A}$;
	\item\ollabel{card:infinite} $A$ Dedekind-sonsuzdur.
\end{enumerate}
\end{thm}

\begin{proof}
\olref[ord-arithmetic][using-addition]{ordinfinitycharacter},
\olref[cardinals][cardsasords]{lem:CardinalsBehaveRight} ve
\olref{naturalsarecardinals} sonuçlarından çıkar.
\end{proof}

Bu, \olref[sfr][][]{part} içinde oldukça gayriresmî kullandığımız bazı
kavramların şu \emph{tanımına} izin verir:

\begin{defn}\ollabel{defnfinite}
$\card{A}$ bir doğal sayıysa, yani $\card{A}\in\omega$ ise $A$'ya
\emph{sonlu}; aksi hâlde \emph{sonsuz} deriz.
\end{defn}
\noindent
Fakat bu tanımın $\ZFC$ arka planında sunulduğuna dikkat edin. Sonuçta her
kümenin bir kardinalitesi olduğunu güvence altına almak için İyi Sıralama'ya
ihtiyaç duyduk. Gerçekten İyi Sıralama olmadan ne sonlu ne de
Dedekind-sonsuz olan bir küme bulunabilir. Bu tür sorulara
\olref[choice][]{chap} içinde döneceğiz. Şimdilik İyi Sıralama'ya
dayanmayı sürdürüyoruz.

Şimdi sonlu kardinal sayılardan sonsuz kardinal sayılara dönelim. İşte iki
temel nokta:

\begin{cor}\ollabel{omegaisacardinal}
$\omega$, en küçük sonsuz kardinal sayıdır.
\end{cor}

\begin{proof}
$\omega$ bir kardinal sayıdır; çünkü $\omega$ Dedekind-sonsuzdur ve bir
$n\in\omega$ için $\cardeq{\omega}{n}$ olsaydı $n$ de Dedekind-sonsuz
olurdu; bu, \olref[z][infinity-again]{naturalnumbersarentinfinite} ile
çelişir. Şimdi tanım gereği $\omega$ en küçük sonsuz kardinal sayıdır.
\end{proof}

\begin{cor}
Her sonsuz kardinal sayı bir limit sıra sayısıdır.
\end{cor}

\begin{proof}
$\alpha$ sonsuz bir ardıl sıra sayısı, dolayısıyla bir $\beta$ için
$\alpha=\beta\ordplus1$ olsun. \olref{finitecardisoequal} gereği $\beta$
de sonsuzdur; böylece
\olref[ord-arithmetic][using-addition]{ordinfinitycharacter} gereği
$\cardeq{\beta}{\beta\ordplus1}$ olur. O zaman
\olref[cardinals][cardsasords]{lem:CardinalsBehaveRight} gereği
$\card{\beta}=\card{\beta\ordplus1}=\card{\alpha}$ olur; dolayısıyla
$\alpha\neq\card{\alpha}$.
\end{proof}

\olref[sfr][siz][enm-alt]{defn:enumerable} kadar erken bir aşamada
!!{enumerable} ve !!{nonenumerable} sonsuz kümeler arasında ayrım
yapabileceğimizi belirtmiştik. Bu tanım doğal olarak şu sonuca yol açar:

\begin{prop}
$A$, !!{enumerable} ancak ve ancak $\card{A}\leq\omega$; $A$,
!!{nonenumerable} ancak ve ancak $\omega<\card{A}$.
\end{prop}

\begin{proof}
Üçlü Karşılaştırma gereği iki iddia denktir; bu yüzden $A$'nın
!!{enumerable} olmasıyla $\card{A}\leq\omega$ olmasının denk olduğunu
ispatlamak yeterlidir. \emph{Sağdan sola}: $\card{A}\leq\omega$ ise
\olref[cardinals][cardsasords]{lem:CardinalsBehaveRight} ve
\olref{omegaisacardinal} gereği $\cardle{A}{\omega}$ olur.
\emph{Soldan sağa}: $A$'nın !!{enumerable} olduğunu varsayın. O zaman
\olref[sfr][siz][enm-alt]{defn:enumerable} gereği üç durum vardır:
\begin{enumerate}
	\item $A=\emptyset$ ise \olref{naturalsarecardinals} ve
	\olref[cardinals][cardsasords]{lem:CardinalsBehaveRight} gereği
	$\card{A}=0\in\omega$ olur.
	\item $\cardeq{n}{A}$ ise \olref{naturalsarecardinals} ve
	\olref[cardinals][cardsasords]{lem:CardinalsBehaveRight} gereği
	$\card{A}=n\in\omega$ olur.
	\item $\cardeq{\omega}{A}$ ise \olref{omegaisacardinal} gereği
	$\card{A}=\omega$ olur.
\end{enumerate}
Böylece her durumda $\card{A}\leq\omega$ olur.
\end{proof}
\noindent
Gerçekten $\omega$'nın özel bir yeri vardır. Sayılabilir birçok sıra sayısı
bulunmasına rağmen:

\begin{cor}
$\omega$, tek !!{enumerable} sonsuz kardinal sayıdır.
\end{cor}

\begin{proof}
$\cardfont{a}$ bir !!a{enumerable} sonsuz kardinal sayı olsun.
$\cardfont{a}$ sonsuz olduğundan $\omega\leq\cardfont{a}$ olur.
$\cardfont{a}$ bir !!a{enumerable} kardinal sayı olduğundan
$\cardfont{a}=\card{\cardfont{a}}\leq\omega$ olur. Üçlü Karşılaştırma
gereği $\cardfont{a}=\omega$ elde edilir.
\end{proof}

Elbette sonsuz sayıda kardinal sayı vardır. Bu nedenle \emph{Kaç kardinal
sayı vardır?} diye sorabiliriz. Aşağıdaki sonuçlar, soruyu yeniden
düşünmemiz gerekebileceğini gösterir.

\begin{prop}\ollabel{unioncardinalscardinal}
$X$'in her elemanı bir kardinal sayıysa $\bigcup X$ de bir kardinal sayıdır.
\end{prop}

\begin{proof}
$\bigcup X$'in bir sıra sayısı olduğunu denetlemek kolaydır.
$\alpha\in\bigcup X$ bir sıra sayısı olsun. O zaman bir $\cardfont{b}$
kardinal sayısı için $\alpha\in\cardfont{b}\in X$ olur.
$\cardfont{b}$ kardinal sayı olduğundan
$\cardless{\alpha}{\cardfont{b}}$ olur. Ayrıca
$\cardfont{b}\subseteq\bigcup X$ olduğundan
$\cardle{\cardfont{b}}{\bigcup X}$ ve dolayısıyla
$\cardneq{\alpha}{\bigcup X}$ olur. Genelleyerek $\bigcup X$'in bir
kardinal sayı olduğu çıkar.
\end{proof}

\begin{thm}\ollabel{lem:NoLargestCardinal}
En büyük kardinal sayı yoktur.
\end{thm}

\begin{proof}
Her $\cardfont{a}$ kardinal sayısı için Cantor Teoremi
(\olref[sfr][siz][car]{thm:cantor}) ile
\olref[cardinals][cardsasords]{lem:CardinalsExist},
$\cardfont{a}<\card{\Pow{\cardfont{a}}}$ olduğunu gerektirir.
\end{proof}

\begin{thm}
Bütün kardinal sayıların kümesi yoktur.
\end{thm}

\begin{proof}
Olmayana ergiyle
$C=\Setabs{\cardfont{a}}{\cardfont{a}\text{ bir kardinal sayıdır}}$
varsayın. \olref{unioncardinalscardinal} gereği $\bigcup C$ bir kardinal
sayıdır. Öyleyse \olref{lem:NoLargestCardinal} gereği
$\cardfont{b}>\bigcup C$ olan bir $\cardfont{b}$ kardinal sayısı vardır.
Tanım gereği $\cardfont{b}\in C$; böylece
$\cardfont{b}\subseteq\bigcup C$ ve dolayısıyla
$\cardfont{b}\leq\bigcup C$ olur. Bu bir çelişkidir.
\end{proof}

Bunu hem Russell Paradoksu'yla hem Burali--Forti ile karşılaştırmalısınız.

